\chapter[\emj{🔢}~Bit Manipulation (Manipulación de Bits)]{\emj{🔢}~Bit Manipulation (Manipulación de Bits)}

\begin{dsaquees}
Los interruptores de luz 💡 de tu casa. 

Cada bit es como un interruptor: 1 significa "encendido" y 0 significa
"apagado". Manipular bits es jugar con esos interruptores usando 
reglas lógicas:
- AND: Solo enciende si AMBOS están encendidos.
- OR: Enciende si AL MENOS UNO está encendido.
- XOR: Solo enciende si EXACTAMENTE UNO está encendido.

Es la forma más profunda y rápida de hablarle a la computadora,
directamente en su lenguaje natural de ceros y unos.
\end{dsaquees}
\begin{dsaotraforma}
Bit Manipulation consiste en aplicar operaciones lógicas (AND, OR, 
XOR, NOT) y de desplazamiento (Shifts) sobre los bits individuales de 
un número entero.

A nivel de hardware, estas operaciones son las más rápidas que un 
procesador puede realizar. Usar bits permite ahorrar muchísima 
memoria y acelerar algoritmos que de otra forma usarían estructuras
más pesadas.
\end{dsaotraforma}
\begin{dsadiagrama}
\begin{verbatim}
Sean A = 5 (101) y B = 3 (011)

AND (\&):  101 \& 011 = 001 (1)
OR  (|):  101 | 011 = 111 (7)
XOR (^):  101 ^ 011 = 110 (6)  <-- ¡Nota: 1^1 = 0!

Left Shift (<<): 5 << 1 = 1010 (10)  (Multiplicar por 2)
Right Shift (>>): 5 >> 1 = 10 (2)    (Dividir entre 2)
\end{verbatim}
\end{dsadiagrama}
\begin{dsacomofunciona}
1. x \textasciicircum{} x = 0: Cualquier número XOR consigo mismo da cero. 
   (Base del problema "Single Number").
2. x \textasciicircum{} 0 = x: Cualquier número XOR cero se queda igual.
3. n \& (n - 1): ¡Truco de Google! Elimina el bit encendido más a la 
   derecha. Si el resultado es 0, entonces `n` era una potencia de 2.
4. i \& 1: Si da 0 es par, si da 1 es impar. (Más rápido que `i \% 2`).
5. n >> 1: Divide entre 2 (truncando).
6. 1 << n: Calcula 2 elevado a la potencia `n`.
\end{dsacomofunciona}
\section*{PSEUDOCÓDIGO}
\hrulefill
\begin{verbatim}
función contarBits(n):
    contador = 0
    MIENTRAS n > 0:
        n = n \& (n - 1) // Eliminar el bit más a la derecha
        contador++
    return contador
\end{verbatim}
\section*{CÓDIGO C++}
\hrulefill
\begin{lstlisting}
#include <iostream>
#include <vector>
using namespace std;

// Problema: Encontrar el número que no se repite
// en un array donde todos los demás aparecen 2 veces.
int findSingleNumber(vector<int>& nums) {
    int res = 0;
    for (int x : nums) {
        res ^= x; // Los duplicados se cancelan: x ^ x = 0
    }
    return res;
}

// Verificar si es potencia de 2
bool isPowerOfTwo(int n) {
    return n > 0 && (n & (n - 1)) == 0;
}

int main() {
    vector<int> data = {4, 1, 2, 1, 2};
    cout << "Número único: " << findSingleNumber(data) << endl; // 4
    
    int val = 16;
    cout << "¿16 es potencia de 2? " << isPowerOfTwo(val) << endl; // 1
    
    return 0;
}
\end{lstlisting}
\begin{dsacomplejidad}
Operación        │ Tiempo  │ Espacio
  ─────────────────┼─────────┼──────────────
  Bitwise (\&,|,\textasciicircum{},\textasciitilde{})│ O(1)    │ O(1)
  Shifts (<<, >>)  │ O(1)    │ O(1)

* Son las operaciones más eficientes posibles en computación.
\end{dsacomplejidad}
\begin{dsaentrevistas}

\end{dsaentrevistas}
\begin{dsaresumen}
• Operaciones a nivel de ceros y unos.
• Súper rápidas O(1).
• AND (\&), OR (|), XOR (\textasciicircum{}), NOT (\textasciitilde{}).
• Shifts: << (x2), >> (/2).
• n \& (n-1) = Truco para potencias de 2.
• XOR cancela duplicados (x \textasciicircum{} x = 0).
• Ideal para problemas de eficiencia extrema.
════════════════════════════════════════════════════════════════
\end{dsaresumen}

